\DocumentMetadata{
    pdfstandard=ua-2,
    % pdfversion=1.7, % Force 1.7 for UA-1 compatibility
    tagging=on,
    testphase={phase-III,table,list},
    lang=en-US,
}

\documentclass{article}

% Fix for the MathML/unicode-math warning 
\tagpdfsetup{math/mathml/luamml/load=false}

% Language and font encoding
\usepackage[british]{babel}

\usepackage{hyperref}      
\hypersetup{
    pdftitle={Semantic Benchmark for PDF/UA}, % Matches \title
    pdfauthor={Floris},
    pdfdisplaydoctitle=true
}
\usepackage{graphicx}      
\usepackage{booktabs}      
\usepackage{xcolor}        
\usepackage{enumitem}      

% Define Title Metadata
\title{Semantic Benchmark for PDF/UA}
\author{Floris}
\date{\today}

\begin{document}

\section{Introduction}
This document is intended as a small semantic benchmark for PDF/UA generation and structure recovery. It contains headings, paragraphs, lists, links, code blocks, a blockquote, a table, a figure, and footnotes.

The main goal is to test whether common authoring pipelines preserve document structure from source to PDF and back again. For a general reference on Markdown, see the \href{https://www.markdownguide.org/basic-syntax/}{Markdown Guide}.

\subsection{Headings and Paragraphs}
A structured document should preserve heading hierarchy, paragraph flow, and inline semantics. This paragraph contains \textbf{bold text}, \textit{italic text}, and \texttt{inline code} to test whether emphasis and code are preserved during conversion.

\subsubsection{Inline Semantics}
Semantic preservation matters because a heading is not just larger text, and a list is not just visually indented text. A machine-readable document should retain these distinctions so that downstream extraction tools can recover them reliably.\footnote{In this benchmark, ``structure'' includes headings, lists, tables, figures, links, reading order, and metadata.}

\section{Lists}
Below is an unordered list with nesting:
\begin{itemize}
    \item First bullet item
    \item Second bullet item
    \begin{itemize}
        \item Nested bullet item
        \item Another nested bullet with a \href{https://www.w3.org/}{link to the W3C}
    \end{itemize}
    \item Third bullet item
\end{itemize}

Below is an ordered list with nesting:
\begin{enumerate}
    \item Prepare the source file.
    \item Export the document to PDF.
    \begin{enumerate}
        \item Enable accessibility tags if the tool supports them.
        \item Record the export settings.
    \end{enumerate}
    \item Validate the resulting PDF.
    \item Compare the extracted output against the source.
\end{enumerate}

\section{Blockquote}
\begin{quote}
    A document can look visually correct while still being structurally broken. Accessibility and machine readability depend on the logical structure being preserved, not only on the page layout.
\end{quote}

\section{Code Block}
The following shell command illustrates a simple conversion workflow:
\begin{verbatim}
pandoc doc-a-semantic-article.md -o doc-a.tex
lualatex doc-a.tex
\end{verbatim}

A second code example in Python:
\begin{verbatim}
from pathlib import Path

source = Path("doc-a-semantic-article.md")
print(f"Processing {source.name}")
\end{verbatim}

\section{Table Example}\label{sec:table-example}
\begin{table}[h]
\centering
\begin{tabular}{lll}
\textbf{Element} & \textbf{Example} & \textbf{Purpose} \\\hline
Heading & '\# Introduction' & Encodes section hierarchy \\
Emphasis & '\textasteriskcentered' & Marks strong importance \\
Link & '[]()' & Preserves navigation and references \\
Footnote & '\^{}1' & Adds supporting notes \\
Code & 'inline' & Distinguishes literal code from text \\
\end{tabular}
\end{table}

\section{Figure Example}
\begin{figure}[ht]
    \centering
    \includegraphics[alt={Coffee and croissant on a bench},width=0.3\textwidth]{../assets/FILE_A_IMG.jpg}
    \caption{Example figure used to test image handling, alt text, and caption preservation.}
    \label{fig:bench}
\end{figure}

\section{Footnotes and Closing}
This sentence includes a footnote for testing purposes.\footnote{A footnote should remain associated with the correct reference point in the main text.} This second sentence includes another footnote to test multiple notes in the same document.\footnote{Some pipelines preserve footnote text but lose the reference structure or numbering semantics.}

In summary, this document is meant to be simple, readable, and rich enough to test whether structure survives across authoring, PDF generation, validation, and extraction workflows


\end{document}